\chapter{绪论}

\section{研究背景与意义}
医学图像检测、分割与量化分析是医学人工智能研究中的核心内容。对于病理图像而言，细胞核检测与实例分割直接关系到组织结构理解、细胞计数、肿瘤微环境分析和后续临床量化建模；对于放射与病理融合分析而言，可靠的区域定位和结构提取又是疗效评估、风险分层和预后预测的重要基础。可以说，前端识别能力的稳定性在很大程度上决定了后续医学分析任务的上限。

然而，这一领域长期存在两个互相交织的难题。其一是监督信号昂贵且稀缺。医学图像尤其是病理图像中的实例级标注，需要专业病理医师或经过长期训练的研究人员完成，时间成本与经济成本都很高。其二是目标域迁移困难。同一疾病在不同中心、不同切片制备流程、不同染色和不同扫描仪条件下，会表现出明显不同的数据分布；即使模型在单一公开数据集上表现良好，一旦进入真实场景，也常常面临性能下降的问题。

近年来，大规模视觉语言预训练模型的兴起为解决上述问题提供了新的可能。一方面，这类模型具备一定的开放词汇识别能力，使得“未见类别识别”和“少样本适配”不再完全依赖人工标注；另一方面，提示学习机制使模型有机会在结构保持相对稳定的前提下，通过较少附加参数实现目标域迁移。这些新进展为医学图像分析带来了新的方法论转向，即从传统强监督闭集学习逐步走向弱监督、无监督和开放场景建模。

本文的研究动机来源于一个持续演化的认识：虽然医学图像强监督标注昂贵，但目标结构、跨模态知识以及少量示例本身蕴含着可被模型利用的先验。换言之，昂贵人工标注并不是唯一的监督来源，合理设计的伪标签、自训练、文本提示、图像提示和结构先验都可能成为替代性学习信号。

沿着这一思路，本文将研究问题组织为一条从“识别入口构建”走向“临床量化解释”的主线：首先，在缺乏目标域检测标注时，如何借助视觉语言大模型的开放词汇能力与自动提示构造建立零样本细胞核检测入口？其次，当零样本病理检测对手工提示工程高度敏感、且密集细胞核易出现漏检与重叠时，如何通过自动提示机制生成更稳定的语义入口，并结合自训练迭代实现目标域适配？再次，当纯文本提示难以稳定表达细粒度医学视觉差异时，如何进一步让图像提示以兼容形式进入视觉语言模型，以增强跨域感知与推理能力？最后，当前端识别能力逐步建立后，如何将其组织为可解释、可统计的数字病理指标，用于疗效评估与预后分析？

\section{国内外研究现状}
\subsection{医学图像检测与分割研究的总体脉络}
医学图像分析的发展大体经历了从手工特征、浅层机器学习到深度学习，再到预训练大模型和提示学习的几个阶段。早期方法主要依赖阈值分割、区域生长、活动轮廓模型、形态学算子以及人工设计纹理特征来完成病灶或组织结构提取。这类方法的优点是可解释性强、实现简单，但它们通常只能适用于成像条件较为稳定、目标形态变化较小的任务。一旦遇到病理染色变化、扫描噪声增强或者目标边界模糊等问题，传统方法往往很难保持稳定性能。

随着卷积神经网络的成功应用，医学图像分析逐渐进入深度学习阶段。以 U-Net 为代表的编码器--解码器结构推动了像素级分割任务的快速发展\cite{ronneberger2015unet}，Mask R-CNN 及其变体则推动了实例级分割的发展\cite{he2017maskrcnn}。对于细胞核检测与分割任务而言，StarDist 的星凸多边形表示和 Cellpose 的通用细胞流场建模进一步提高了边界恢复能力和复杂背景下的识别能力\cite{schmidt2018stardist,stringer2021cellpose}。nnU-Net 则通过数据指纹驱动的自动配置机制，说明稳健的预处理、网络结构和训练策略对跨任务医学分割同样关键\cite{isensee2021nnunet}。然而，这一阶段的方法几乎都建立在“存在足够强监督标注”的前提之上，也因此把问题转移到了数据构建成本上。医学数据的标注瓶颈并没有因为模型变强而消失，反而由于模型更依赖大规模训练而变得更加突出。

近几年，视觉语言预训练和基础模型的发展又把医学图像分析推入新的阶段。研究重点开始从“如何在固定数据集上逼近最优精度”转向“如何减少对强监督标注的依赖”“如何在开放类别和跨域场景中泛化”以及“如何利用文本、图像示例和外部知识协助推理”。本文所关注的弱监督、零样本、图像提示和数字病理分析问题，本质上都处在这一新阶段的演进链条中。

\subsection{弱监督与无监督医学图像分割}
\subsubsection{弱监督分割的基本思想}
弱监督分割的基本目标是利用比像素级标注更廉价的监督信息近似恢复目标区域。常见弱监督信号包括图像级标签、点标注、框标注、涂鸦标注以及部分伪标签。不同监督形式对应不同的信息上限：图像级标签最廉价，但只告诉模型某类目标是否存在；点和框能提供更明确的空间约束，但边界信息仍然缺失；涂鸦比点和框更接近真实轮廓，但人工成本也更高。如何在“监督强度”和“标注成本”之间取得平衡，是弱监督分割长期关注的问题。

在自然图像领域，类激活图、区域提案、条件随机场后处理和自训练策略构成了图像级弱监督分割的主流技术路线。即先通过分类网络获得目标的粗激活分布，再利用后处理或迭代优化方式把粗热图转化为较完整的掩码。在医学图像场景中，这一思路同样具有启发性，但直接照搬并不总是有效。医学对象尤其是细胞核目标具有尺度小、密度高、粘连严重等特点，使得自然图像中的弱监督前景恢复策略很难直接实现高质量实例边界提取。

\subsubsection{病理图像中的特殊困难}
病理图像弱监督分割面临的困难具有鲜明的结构性。首先，细胞核等目标往往尺寸小、数量多，相互之间距离很近，一旦伪标签稍微膨胀，就会把多个目标错误地连在一起。其次，病理图像中的颜色和纹理受染色和扫描条件影响显著，不同样本之间的成像风格变化容易导致弱监督定位结果不稳定。再次，病理图像中存在大量与目标外观接近的背景结构，这些结构往往会在弱监督阶段被错误激活，造成伪标签污染。

因此，病理图像中的弱监督分割方法通常不能只停留在“恢复前景区域”层面，而必须进一步思考如何形成实例级结构感知。这也是为什么细胞核分割工作中常常会显式引入边界分支、距离变换分支或中心点约束等结构先验。类似地，本文后续针对密集细胞核场景中的漏检与重叠问题，也引入了基于伪标签迭代的自训练/蒸馏机制，以在弱监督信号有限的条件下逐步提升结构化定位质量。

\subsubsection{伪标签与自训练}
伪标签和自训练在弱监督与半监督分割中占据重要位置。其基本思想是先利用现有模型生成高置信预测，再把这些预测当作新监督信号继续训练模型。成功的关键在于：伪标签既不能过于保守，否则覆盖率不足；也不能过于激进，否则错误会在迭代中不断放大。在医学图像中，这个平衡更难把握，因为目标形态复杂且域偏移明显。

已有研究证明，自训练在医学图像中并非简单的“多跑几轮”策略，而需要与结构先验、分支解耦或多任务约束结合。对于本文的研究主线而言，伪标签与自训练并不是一个孤立技术点，而是一种贯穿多篇工作的核心思想：从弱监督分割中的伪标签迭代优化，到零样本检测后的目标域自适配（自训练/蒸馏），再到图像提示学习中的少样本引导，都能看到“先构造一个可接受起点，再通过数据反馈逐步提升”的共通逻辑。

\subsection{细胞核检测与实例分割研究}
\subsubsection{检测、分割与联合建模}
细胞核分析任务通常分为三类：核检测、语义分割和实例分割。核检测强调目标位置，适合计数和密度估计；语义分割强调前景区域覆盖，适合估计总体核面积；实例分割则要求区分每个独立核目标，最适合后续的形态统计和组织微环境分析。三类任务并非彼此割裂，而是存在显著内在联系。一个好的实例分割模型往往隐含具有检测能力，一个好的检测器也通常能为实例学习提供较稳定候选区域。

在病理图像中，许多方法试图将检测与分割联合建模。例如，有的工作先进行核中心点检测，再以中心点为引导恢复边界；有的工作同时预测前景图、边界图和距离图，以实现实例拆分；还有的工作直接通过 proposal-based 结构完成实例掩码预测。不同方法的差异在于对实例结构的建模方式不同，但它们都说明一件事：对于高密度细胞核任务，仅靠单一前景预测往往不够，必须引入额外结构信息。

\subsubsection{跨数据集与跨域问题}
细胞核数据集之间的差异很大。PanNuke、MoNuSAC 和 Lizard 等公开基准覆盖了不同器官、组织类型与核类别，为实例分割和分类研究提供了重要基础\cite{gamper2019pannuke,verma2021monusac,graham2021lizard}。2018 Data Science Bowl 的跨成像实验核分割评测进一步揭示了通用核分割模型在成像条件变化下的泛化难点\cite{caicedo2019datasciencebowl}。然而，不同组织来源、不同切片质量、不同染色方案和不同扫描仪仍会导致目标外观变化。一个在单一公开 benchmark 上表现优异的模型，进入真实跨域场景后往往会遇到性能下降。因此，单纯追求单数据集分数并不足以说明方法的实用性。如何在有限目标域支持样本甚至无支持样本条件下增强泛化能力，已经成为细胞核分析研究的重要方向。

这一点与本文后续关于跨域评测（domain shift）、少样本支持提示（few-shot prompting）以及“冻结主干、仅通过提示适配”的直测设定密切相关。作者之所以反复强调源域训练、目标域少量支持样本提示与不进行目标域大规模微调，是因为细胞核分析任务的实际挑战并不仅仅是“在一个数据集上做得更高”，而是“面对新域时还能否工作”。从学位论文角度看，这种问题意识也使得全文研究比单纯追求基准分数更具完整性。

\subsection{视觉语言预训练模型及其医学迁移}
\subsubsection{视觉语言预训练的基本范式}
视觉语言预训练模型通过大规模图文对构建跨模态共享语义空间。CLIP 通过图像编码器和文本编码器的对比学习训练获得强分类迁移能力，BLIP 在此基础上进一步增强生成与理解能力，而 GLIP 则尝试把语言提示直接融入目标检测框架，使检测器具备开放词汇识别能力\cite{radford2021learning,li2022glip}。RegionCLIP 进一步把图文预训练推进到区域级语义对齐，为开放词汇定位提供了更细粒度的表示基础\cite{zhong2022regionclip}。OWL-ViT 通过简洁的视觉 Transformer 结构实现文本条件和图像条件的开放词汇检测，分层视觉语言蒸馏方法则探索了将预训练语义迁移到单阶段检测器的路径\cite{minderer2022owlvit,ma2022hierkd}。这一系列模型的重要意义在于，它们使图像理解不再局限于固定类别表，而可以借助文本描述扩展到未见对象。

在自然图像领域，开放词汇检测的价值已经得到广泛验证。但在医学图像领域，这一能力并不能直接迁移。首先，医学图像的许多概念在大规模通用图文语料中本来就稀缺；其次，医学对象往往依赖局部结构模式而非简单语义标签；再次，医学任务常常需要非常稳定的定位能力，而非仅仅粗粒度识别。因此，视觉语言模型迁移到医学场景时，需要配合新的提示机制、目标域自适应机制和可能的结构先验。近年来，PLIP、CONCH、UNI 与 Virchow 等计算病理基础模型分别从病理图文对齐、大规模组织表征和临床级泛化等角度验证了领域预训练的价值\cite{huang2023plip,lu2024conch,chen2024uni,vorontsov2024virchow}。CTransPath、HIPT 与 Prov-GigaPath 又分别从局部--全局对比学习、WSI 层级建模和超长上下文预训练角度拓展了计算病理表征学习的尺度\cite{wang2022ctranspath,chen2022hipt,xu2024gigapath}。

\subsubsection{开放词汇与零样本能力的局限}
尽管视觉语言模型具备开放词汇能力，但在医学任务中，单靠文本类名通常不足以支撑高质量检测。例如，同样是“细胞核”，不同病理环境下的核在大小、形状、染色情况和周边背景方面可能完全不同。简单提示无法表达这种差异，甚至会导致模型偏向自然图像语义空间中的错误对象。因此，零样本医学检测如果要真正成立，必须在提示构造和目标域适配上做更细致设计。

本文第三章正是在这一局限之上提出自动属性提示构造。其出发点不是否认视觉语言模型的价值，而是承认“医学类别语义不足”的现实，然后试图以属性、描述和自训练的方式弥补这一不足。从相关工作回顾角度看，这使得本文不是简单使用现成大模型，而是在其基础上提出新的迁移方法。

\subsection{提示学习的发展与图像提示研究}
\subsubsection{文本提示学习}
提示学习最早主要体现为对文本模板的设计。研究者通过改变 prompt phrasing 来影响模型对类别的理解，例如“a photo of a dog”和“this is a dog”在 CLIP 类模型中可能带来不同效果。随着研究发展，更多工作开始引入可学习 prompt token，以代替人工设计模板，并通过条件提示使提示表示随输入图像自适应变化\cite{zhou2022cocoop}。对医学图像而言，文本提示学习具有明显吸引力，因为它不需要大规模改动模型结构，却可以借助预训练语义进行目标域适配。

不过，医学场景中的文本提示存在两个根本局限。第一，语言本身未必能精确表达细粒度视觉差异。第二，许多医学任务对形态和组织结构极为敏感，而这些信息天然更容易通过视觉示例表达，而非通过词语表达。因此，仅依赖文本 prompt 往往不足以满足复杂医学场景需求。

\subsubsection{图像提示与多模态提示}
图像提示的核心思想是让支持图像本身参与推理，以补充纯文本提示不足。已有工作在通用目标检测、分割和 few-shot 学习中表明，少量支持样本可以显著改善模型对新类别的感知能力；通用分割基础模型及其医学适配版本也展示了视觉提示驱动分割的可迁移性\cite{kirillov2023sam,ma2024medsam}。但如何让图像提示与原有视觉语言预训练结构兼容，仍是一个开放问题。若处理不当，图像提示可能破坏原有图文对齐能力，导致模型难以稳定收敛。

在医学场景下，图像提示的潜力更加明显。一个病理专家往往只需要看几张代表性示例就能理解某类目标的大致外观，而图像提示恰好模拟了这种“示例式知识传递”过程。本文第四章的视觉文本化工作，就是对这一思路的系统展开：既不放弃文本提示提供的语义框架，又引入图像提示补充细粒度结构信息。

\subsection{数字病理分析与可解释医学人工智能}
\subsubsection{从识别到量化}
数字病理研究正在从“自动识别病灶”向“自动量化病理表型”演化。弱监督全切片学习已经能够在大规模 WSI 上完成临床级分类与数据高效建模\cite{campanella2019clinical,lu2021clam}。注意力 MIL、DSMIL 与 TransMIL 分别从可解释聚合、实例关系建模和空间相关性建模角度完善了仅使用切片级标签的学习范式\cite{ilse2018attentionmil,li2021dsmil,shao2021transmil}。CAMELYON16 与 PANDA 等多中心评测也表明，算法进入临床病理场景时必须同时关注准确性、外部泛化和评价规范\cite{bejnordi2017camelyon,bulten2022panda}；肺癌病理图像则被证明可用于组织学分类和分子改变预测\cite{coudray2018lung}。前者关注的是目标是否被正确找出，后者关注的是目标识别后能否形成对医生有意义的统计指标。例如，肿瘤残留比例、免疫细胞密度、肿瘤边界结构变化以及细胞空间共定位模式，都属于后者关心的内容。

这种演化意味着，前端检测和分割结果的意义不再只体现为 benchmark 分数，而是要进一步考虑结果是否足以支持后续统计分析与临床解释。因此，一个看似精度很高但结构不稳定的模型，未必真的适合数字病理分析；相反，一个在临床关键指标上更稳定、更可解释的模型，可能更有实际应用价值。

\subsubsection{临床结局建模}
在肿瘤研究中，治疗后病理切片提供了非常丰富的预后信息。除了传统的病理缓解程度评估，肿瘤周围免疫细胞的数量、密度和空间分布也逐渐被认为与治疗反应和生存结局有关。AI 模型若能稳定识别并量化这些指标，就能为更细致的风险分层和个体化治疗提供辅助。

本文第五章的肺鳞癌工作正是建立在这一思想之上。它使得前几章偏方法论的研究，有了一个自然的应用收束点，也让整篇学位论文的价值不局限于算法创新，而延伸到临床解释和数字病理转化。

\section{研究内容与主要贡献}
围绕上述主线，本文主要开展以下四方面研究。

第一，研究视觉语言大模型在细胞核检测中的迁移问题，提出基于自动属性提示与自训练的零样本细胞核检测方法。该工作解决的是“缺乏目标域检测标注时如何构造可辨识的医学提示，并将开放词汇能力转化为目标域可用的检测性能”。

第二，针对零样本检测对提示工程的敏感性以及密集细胞核场景下的漏检、重叠问题，提出自动提示框架 AttriPrompter：从无标注病理图像中自动生成形状/颜色属性，并通过同义与“程度”增强扩展描述粒度，进一步按图文相关性对提示序列排序以获得更可靠的初始定位；在此基础上，引入自训练知识蒸馏迭代优化，以伪标签驱动目标域适配并提升无监督检测性能。该工作解决的是“如何在无需人工提示模板与检测标注的条件下，自动构造高质量提示并稳定提升零样本检测效果”。

第三，研究图像提示与文本提示的统一建模与注入机制，提出视觉文本化方法，通过将少量支持图像映射到语言空间兼容的提示表示，增强模型对细粒度差异与跨域风格变化的感知能力。该工作解决的是“当文本描述不足时，如何让图像示例以提示的形式直接参与推理”。

第四，将上述方法体系进一步应用于肺鳞癌新辅助免疫化疗后的数字病理分析，构建面向残余肿瘤负荷与肿瘤周围免疫细胞密度的量化框架，并通过统计分析探索其与疗效及预后之间的关系。该工作解决的是“如何将检测、分割结果转化为临床可解释指标并形成分析闭环”。

本文的主要贡献体现在以下四个方面。其一，提出了基于自动属性提示与自训练的零样本细胞核检测方法，为医学无标注检测提供了可扩展的实现路径。其二，提出了自动提示框架 AttriPrompter，通过“属性生成—增强—排序”自动构造高质量提示序列，并结合自训练知识蒸馏提升密集细胞核场景下的零样本检测稳定性与目标域适配能力。其三，提出了视觉文本化方法，将图像提示映射到语言空间，从而增强模型在跨域与细粒度目标场景中的感知能力。其四，构建了面向肺鳞癌疗效评估的数字病理分析框架，将前端识别结果进一步组织为具有临床解释意义的量化指标并开展统计分析。




\section{论文组织结构}
本文的技术路线可以概括为“从提示驱动的识别能力构建到临床导向的可解释量化”的逐步推进：前期工作面向无检测标注场景，重点解决自动提示构造与目标域自适配问题；中期工作进一步面向未见类别与跨域分布偏移，重点解决结构先验注入与提示融合问题；后期工作将上述识别能力整合到临床导向的数字病理链条中，强调可解释指标构建与统计验证。

从建模角度看，本文的方法体系有三个共同特点。第一，强调尽量减少对昂贵强监督标注的依赖。第二，强调最大化利用视觉语言预训练带来的跨模态先验，并通过提示机制进行轻量适配。第三，强调输出结果的可解释性，使方法不仅服务于 benchmark 指标，也服务于真实数字病理分析需求。


图~\ref{fig:thesis_research_framework} 按照“研究内容—研究思路—对应章节—解决问题”的逻辑概括全文。第二章首先以自动属性提示构造无标注场景下的语义入口；第三章在此基础上引入提示增强、预测反馈与自训练，使零样本识别由一次推理扩展为可迭代优化过程；第四章进一步把支持图像转化为语言空间兼容的视觉 token，以应对文本难以完整表达的外观差异和跨域偏移；第五章则把前述检测、分割能力嵌入肺鳞癌全切片分析流程，形成从算法输出到临床病理指标的转化链条。

\begin{figure}[htbp]
  \centering
  \resizebox{\textwidth}{!}{%
  \begin{tikzpicture}[
    x=1cm,y=1cm,
    every node/.style={font=\small,align=center},
    titlebox/.style={draw=black!70,line width=0.55pt,fill=white,minimum width=9.6cm,minimum height=0.72cm},
    introbox/.style={draw=black!55,line width=0.45pt,fill=gray!6,minimum width=8.0cm,minimum height=0.68cm},
    contentbox/.style={draw=black!50,line width=0.4pt,fill=blue!3,text width=3.05cm,minimum height=1.05cm},
    ideabox/.style={draw=black!48,dashed,line width=0.4pt,fill=white,text width=3.05cm,minimum height=1.05cm},
    chapterbox/.style={draw=black!60,line width=0.45pt,fill=gray!10,text width=3.05cm,minimum height=1.05cm},
    problembox/.style={draw=black!50,line width=0.4pt,fill=orange!4,text width=3.05cm,minimum height=0.96cm},
    sidebox/.style={draw=black!35,line width=0.35pt,fill=gray!4,text width=1.25cm,minimum height=0.72cm},
    arrow/.style={-{Stealth[length=1.8mm]},line width=0.45pt,draw=black!65}
  ]
    \node[titlebox] (title) at (8.4,10.3) {基于视觉语言预训练与提示学习的无监督医学图像检测、分割研究};
    \node[introbox] (intro) at (8.4,9.15) {第一章：绪论——研究背景、国内外现状、研究内容与论文组织};
    \node[sidebox] at (0.75,7.65) {研究内容};
    \node[sidebox] at (0.75,6.10) {研究思路};
    \node[sidebox] at (0.75,4.45) {对应章节};
    \node[sidebox] at (0.75,2.85) {解决问题};
    \node[contentbox] (c2) at (3.0,7.65) {自动属性提示与零样本细胞核检测};
    \node[contentbox] (c3) at (6.6,7.65) {提示增强、实例反馈与迭代自训练};
    \node[contentbox] (c4) at (10.2,7.65) {视觉文本化的图像提示学习};
    \node[contentbox] (c5) at (13.8,7.65) {肺鳞癌数字病理量化与预后分析};
    \node[ideabox] (i2) at (3.0,6.10) {属性语义生成 + 伪标签自训练};
    \node[ideabox] (i3) at (6.6,6.10) {提示排序 + 预测反馈 + 知识蒸馏};
    \node[ideabox] (i4) at (10.2,6.10) {支持图像映射为语言空间视觉 token};
    \node[ideabox] (i5) at (13.8,6.10) {多尺度 WSI 分析 + 细胞空间统计};
    \node[chapterbox] (w2) at (3.0,4.45) {第二章\\自动属性提示驱动的零样本细胞核检测};
    \node[chapterbox] (w3) at (6.6,4.45) {第三章\\反馈式属性提示与无监督迭代优化};
    \node[chapterbox] (w4) at (10.2,4.45) {第四章\\视觉文本化图像提示的跨域检测};
    \node[chapterbox] (w5) at (13.8,4.45) {第五章\\肺鳞癌疗效与空间免疫数字病理分析};
    \node[problembox] (p2) at (3.0,2.85) {目标域强监督缺失\\医学提示难以构造};
    \node[problembox] (p3) at (6.6,2.85) {提示敏感、密集目标漏检\\与预测不稳定};
    \node[problembox] (p4) at (10.2,2.85) {文本粒度不足\\与跨域外观偏移};
    \node[problembox] (p5) at (13.8,2.85) {识别结果难以转化为\\临床可解释指标};
    \node[introbox] (end) at (8.4,1.20) {第六章：总结与展望};
    \draw[arrow] (title) -- (intro);
    \draw[arrow] (intro.south) -- ++(0,-0.32) -| (c2.north);
    \draw[arrow] (intro.south) -- ++(0,-0.32) -| (c3.north);
    \draw[arrow] (intro.south) -- ++(0,-0.32) -| (c4.north);
    \draw[arrow] (intro.south) -- ++(0,-0.32) -| (c5.north);
    \foreach \n in {2,3,4,5} {
      \draw[arrow] (c\n) -- (i\n);
      \draw[arrow] (i\n) -- (w\n);
      \draw[arrow] (w\n) -- (p\n);
    }
    \draw[arrow] (w2.east) -- (w3.west);
    \draw[arrow] (w3.east) -- (w4.west);
    \draw[arrow] (w4.east) -- (w5.west);
    \draw[arrow] (8.4,2.15) -- (end);
  \end{tikzpicture}}
  \caption{本文研究内容与组织结构}
  \label{fig:thesis_research_framework}
\end{figure}

从章节关系看，第二章和第三章共同研究低标注成本条件下的细胞核识别，但两者侧重点不同：第二章关注从视觉语言先验中建立可用的零样本检测入口，第三章强调利用实例特征、中文提示、自训练和分割头等补充机制形成反馈式优化。第四章把研究对象由核图像扩展到通用与医学跨域目标检测，突出图像提示对未见类别的适配能力。第五章进一步转向临床应用，围绕残余活性肿瘤、肿瘤周围细胞及其空间分布建立可解释量化体系。

全文共分六章。第一章为绪论，介绍研究背景与意义、国内外研究现状、主要研究内容和论文组织结构。第二章介绍基于自动属性提示与自训练的零样本细胞核检测方法。第三章介绍自动提示框架 AttriPrompter 及其自训练知识蒸馏迭代优化方法。第四章介绍视觉文本化框架及其图像提示驱动的检测与分割方法。第五章介绍面向肺鳞癌疗效评估的数字病理量化分析框架。第六章总结全文并讨论后续研究方向。
